\chapter{Référence de l'API Lua}
\section{Modèle d'exécution}
app.lua renvoie une table M. M.loop() est obligatoire et s'exécute de façon répétée ; M.init(), M.event(event\_type) et M.cleanup() sont facultatifs. Sans loop, l'application est rejetée. En cas d'erreur d'exécution, le micrologiciel peut basculer vers une application de secours.

Le script dans sensor/*.lua renvoie une table sensor avec sensor.get(index), qui renvoie une valeur numérique. Le micrologiciel expose sensor globalement. apppath est le chemin du dossier de l'application courante. Charger les modules avec require("module") ; pcall(require, "module") convient aux modules facultatifs.
\section{Ressources et affichage}
loadimg et loadfont renvoient un identifiant ou nil ; free libère les images et polices chargées. fill remplit l'écran entier. update transmet l'image puis efface l'écran en noir. drawpolygon demande un nombre pair de coordonnées, au moins quatre éléments. Appeler update une fois par image et libérer les ressources dans cleanup.
\section{Autres modules}
system contrôle délai, rétroéclairage, résistance de tirage V2, sortie 5 V et veille, et lit boutons, ADC et heure. onewire recherche les DS18B20 et lit leur température en Celsius. nvs utilise l'espace NVS app ; get renvoie une chaîne ou nil et push un booléen. ble initialise BLE et fournit régime, vitesse et rapport.
\section{display}
\begin{itemize}
\item display.loadimg(path) -> img\_handle|nil
\item display.loadfont(path) -> font\_handle|nil
\item display.free(handle)
\item display.fill(color)
\item display.update()
\item display.drawimg(img, x, y)
\item display.string(font, x, y, text, color)
\item display.fillrect(x, y, width, height, color)
\item display.drawline(x1, y1, x2, y2, color)
\item display.drawcircle(x, y, r, color)
\item display.drawfillcircle(x, y, r, color)
\item display.drawroundrect(x1, y1, x2, y2, r, color)
\item display.drawarrow(x0, y0, x1, y1, w, color)
\item display.drawfillarrow(x0, y0, x1, y1, w, color)
\item display.drawpolygon(\{x1, y1, x2, y2, ...\}, color)
\item display.RED
\item display.GREEN
\item display.BLUE
\item display.WHITE
\item display.BLACK
\item display.YELLOW
\item display.CYAN
\item display.PURPLE
\item display.GRAY
\end{itemize}
\section{system}
\begin{itemize}
\item system.delay(ms)
\item system.backlight(level)
\item system.V2\_pullup(enable)
\item system.out5V(enable)
\item system.keys() -> int
\item system.adc\_ch(channel) -> number
\item system.get\_time() -> hour, minute, second
\item system.config\_sleep(enable)
\end{itemize}
\section{onewire}
\begin{itemize}
\item onewire.scan18b20() -> count
\item onewire.read18b20(index) -> temperature\_c
\end{itemize}
\section{nvs}
\begin{itemize}
\item nvs.get(key) -> string|nil
\item nvs.push(key, value\_string) -> boolean
\end{itemize}
\section{ble}
\begin{itemize}
\item ble.init() -> true
\item ble.get\_telemetry() -> rpm, speed, gear
\end{itemize}
\section{Exemples et conseils}
L'exemple de voltmètre utilise des coordonnées dans les 160\,\texttimes\,80 pixels et borne la largeur de la barre. L'exemple de capteur montre lecture ADC et lissage, pas une conversion NTC complète : remplacer la transformation illustrative par l'étalonnage réel. Garder sensor.get rapide ; stocker les étalonnages avec nvs et convertir les chaînes en nombres si nécessaire. Construire les chemins depuis apppath. Rechercher les DS18B20 une fois puis les lire dans la boucle.
\begin{verbatim}
local M = {}
local font = display.loadfont(apppath .. "/../fonts/H24XB.FNT")

local function read_voltage()
return sensor.get(1)
end

function M.loop()
local v = read_voltage()
display.fill(display.BLACK)
display.string(font, 10, 30,
    string.format("%2.1f V", v), display.WHITE)
display.fillrect(10, 60,
    math.max(0, math.min(140, math.floor((v - 9.5) * 12))),
    8, display.GREEN)
display.update()
end

function M.cleanup()
display.free(font)
end

return M
\end{verbatim}
\begin{verbatim}
sensor = {}

local filt = 0.0

function sensor.get(i)
i = i or 0
system.V2_pullup(1)

local raw = system.adc_ch(i)
filt = filt + (raw - filt) / 10.0

local value = filt * 1.0
return value
end

return sensor
\end{verbatim}
